% Phase 2 drafts (T1--T4). NOT included by any build: no \input refers to
% this file, and it sits outside the paper/*.tex glob. Notation follows
% paper/transport_theory.tex. Proof-status tags:
%   [checked]      every step was checked line by line in this session;
%   [not checked]  the step relies on a statement whose proof was not
%                  re-checked in this session.

\section*{Draft T1: generalized-linear and logistic bandit literature}

% Sources opened in this session (arXiv PDFs, text extracted with pypdf):
%   Li, Lu, Zhou, arXiv:1703.00048v2; first page reads "Proceedings of the
%     34th International Conference on Machine Learning, Sydney, Australia,
%     PMLR 70, 2017".
%   Faury, Abeille, Calauzenes, Fercoq, arXiv:2002.07530v2 (8 Jun 2020); the
%     arXiv text carries no venue line and the arXiv API returns no journal_ref.
%   Abeille, Faury, Calauzenes, arXiv:2010.12642v2; first page reads
%     "Proceedings of the 24th International Conference on Artificial
%     Intelligence and Statistics (AISTATS) 2021, San Diego, California, USA.
%     PMLR: Volume 130".
% NOT opened: Filippi, Cappe, Garivier, Szepesvari (2010). Every NeurIPS
%   proceedings, HAL and author-page URL tried was refused by the proxy.
% NOT verified: no entry was checked against DBLP or a proceedings page;
%   dblp.org and proceedings.mlr.press were refused (proxy HTTP 403;
%   WebFetch domain refusal). Page ranges are therefore omitted below.
%
% Draft bibliography entries (fields limited to what the opened PDFs state):
% @inproceedings{li2017provably,
%   title     = {Provably Optimal Algorithms for Generalized Linear Contextual Bandits},
%   author    = {Li, Lihong and Lu, Yu and Zhou, Dengyong},
%   booktitle = {Proceedings of the 34th International Conference on Machine Learning},
%   series    = {Proceedings of Machine Learning Research}, volume = {70}, year = {2017}}
% @misc{faury2020improved,
%   title  = {Improved Optimistic Algorithms for Logistic Bandits},
%   author = {Faury, Louis and Abeille, Marc and Calauz{\`e}nes, Cl{\'e}ment and Fercoq, Olivier},
%   note   = {arXiv:2002.07530v2}, year = {2020}}   % venue unverified
% @inproceedings{abeille2021instance,
%   title     = {Instance-Wise Minimax-Optimal Algorithms for Logistic Bandits},
%   author    = {Abeille, Marc and Faury, Louis and Calauz{\`e}nes, Cl{\'e}ment},
%   booktitle = {Proceedings of the 24th International Conference on Artificial Intelligence and Statistics},
%   series    = {Proceedings of Machine Learning Research}, volume = {130}, year = {2021}}
% filippi2010parametric: not drafted; primary source not opened.

\paragraph{Target model.} Rewards are $r_t=\mu_{\btheta^\circ,W}(z_t)+\eta_t$ with
$\mu_{\btheta,W}(z)=\sqrt W\tanh(\phi(z)^\top\btheta/\sqrt W)$, Gaussian
$\eta_t\sim N(0,\sigma^2)$, $\|\phi\|_2\le B_\phi$, $\|\btheta^\circ\|_2\le R$.
This is the single-index (generalized-linear) form $E[r\mid z]=g(\phi(z)^\top\btheta^\circ)$
with known, strictly increasing link $g(u)=\sqrt W\tanh(u/\sqrt W)$,
$0<\dot g\le1$, $|\ddot g|\le c_h/\sqrt W$ with $c_h=4/(3\sqrt3)$ at $B_\phi=1$,
and $|\ddot g|\le(2/\sqrt W)\dot g$. The paper fits it with squared loss and an
arbitrary predictable parameter; it does not compute a maximum-likelihood or
quasi-likelihood estimate.

\begin{description}
\item[Li, Lu and Zhou (2017).]
\emph{Noise:} $Y_t=\mu(X_t^\top\theta^*)+\epsilon_t$ with
$E[\epsilon_t\mid\mathcal F_{t-1}]=0$ and conditional $\sigma$-sub-Gaussian
$\epsilon_t$ (their Eq.~(2)); the problem setting also states rewards
$y_{t,a}\in[0,1]$. Contexts are drawn i.i.d.\ from $\nu$ with $\|x_{t,a}\|\le1$
and $\lambda_{\min}(E[\frac1K\sum_a x_{t,a}x_{t,a}^\top])\ge\sigma_0^2>0$.
\emph{Link:} fixed and strictly increasing; Assumption~1:
$\kappa:=\inf_{\|x\|\le1,\|\theta-\theta^*\|\le1}\dot\mu(x^\top\theta)>0$;
Assumption~2: $\mu$ twice differentiable with $\dot\mu\le L_\mu$,
$|\ddot\mu|\le M_\mu$. \emph{Estimator:} the root of
$\sum_{i<t}(Y_i-\mu(X_i^\top\theta))X_i=0$ (their Eq.~(6)).
\emph{Bounds:} UCB-GLM (Theorem~2), with probability at least $1-2\delta$,
$R_T\le\tau+\frac{2L_\mu\sigma d}{\kappa}\log(\frac{T}{d\delta})\sqrt T$ for
$\tau=C\sigma_0^{-2}(d+\log(1/\delta))$; SupCB-GLM (Theorem~3), with probability
at least $1-\delta$ and $T\ge T_0$, a bound $45(\sigma L_\mu/\kappa)\sqrt{dT}$
times logarithmic factors in $T$, $K$, $d$ and $1/\delta$ (the extent of a
radical in that display is ambiguous in the extracted text and is not
transcribed here), stated to give
$E[R_T]=O((\log T)^{1.5}\sqrt{dT\log K})$ for $\delta=1/T$.
\emph{Coverage of Gaussian scaled tanh with squared loss: not covered as stated.}
The link conditions hold ($L_\mu=1$, $M_\mu=c_h/\sqrt W$, and with
$\|\btheta^\circ\|=1$, $\kappa=\operatorname{sech}^2(2/\sqrt W)$: $0.786$,
$0.0707$, $0.00134$ at $W=16,1,0.25$), and Gaussian noise is sub-Gaussian.
Three conditions fail or differ. (i) The setting states $y\in[0,1]$, and the
proof of Theorem~2 bounds the initialization regret using $\mu\in[0,1]$;
Gaussian rewards are unbounded. (ii) In the study's 29 coordinates the 80
feature vectors have rank 25, so $\sigma_0=0$; after restriction to their
span, $\sigma_0^2=1/60$ (computed from \texttt{feature\_table}). (iii) Their
estimator solves Eq.~(6), not squared loss, and their guarantee is for their
algorithm, not for the paper's corrected center. Whether the appendix proof of
their Lemma~3 uses bounded rewards was not checked.
\item[Faury, Abeille, Calauz\`enes and Fercoq (2020).]
\emph{Noise:} $r_{t+1}\in\{0,1\}$ with $P(r_{t+1}=1\mid x_t,\mathcal F_t)=\mu(x_t^\top\theta_*)$.
\emph{Link:} the sigmoid $\mu(x)=(1+e^{-x})^{-1}$, $L=M=1/4$; $\theta_*$ in a
compact $\Theta$ with known $S$; $\|x\|_2\le1$;
$\kappa:=\sup_{x\in\mathcal X,\theta\in\Theta}1/\dot\mu(x^\top\theta)$. Their
concentration result (Theorem~1) requires martingale increments with
$|\varepsilon_{t+1}|\le1$ almost surely. \emph{Bounds:} Logistic-UCB-1
(Theorem~2), w.p.\ $\ge1-\delta$, $R_T\le C_1L\sqrt\kappa\,\gamma_T(\delta)\sqrt T$,
$O(\sqrt\kappa\,d\sqrt T\log T)$ for $\lambda=d\log T$; Logistic-UCB-2
(Theorem~3), $R_T\le C_2\gamma_T(\delta)\sqrt T+C_3\gamma_T(\delta)^2\kappa$, with
$\gamma_t(\delta)$ defined after their Eq.~(4).
\emph{Coverage: not covered.} The reward model is Bernoulli with the sigmoid
link, and Theorem~1 needs bounded increments; Gaussian noise violates both.
The affine identity $\sqrt W\tanh(u/\sqrt W)=\sqrt W(2\mu(2u/\sqrt W)-1)$ does
not change the noise model.
\item[Abeille, Faury and Calauz\`enes (2021).]
\emph{Noise:} Bernoulli rewards with mean $\mu(x_t^\top\theta_\star)$,
$\mu$ the logistic function, using $|\ddot\mu|\le\dot\mu$. \emph{Assumptions:}
$\|x\|\le1$; $\|\theta_\star\|\le S$ known. \emph{Bounds:} OFULog (Theorem~1):
$\mathrm{Regret}\le R^{\rm perm}+R^{\rm trans}$ with, with high probability and
up to logarithmic factors ($\lesssim_T$), $R^{\rm perm}\lesssim d\sqrt{T/\kappa_\star}$
and $R^{\rm trans}\lesssim\kappa_{\mathcal X}d^2\wedge(d^2+\sum_t\mathbf 1(x_t\in\mathcal X_-))$,
where $\kappa_\star=1/\dot\mu(x_\star^\top\theta_\star)$ and
$\kappa_{\mathcal X}=\max_x1/\dot\mu(x^\top\theta_\star)$. Theorem~2: for
$\mathcal X=S^d(0,1)$ and $T\ge d^2\kappa_\star$, a local minimax lower bound
$\Omega(d\sqrt{T/\kappa_\star})$. \emph{Coverage: not covered}, for the same
reasons as Faury et al.: Bernoulli noise, logistic link.
\item[Filippi, Capp\'e, Garivier and Szepesv\'ari (2010).]
\emph{Primary source not opened; no statement is drafted from it.} Secondary
statements in Li et al.\ (2017, pp.~2--5), not verified against the source:
GLM-UCB attains $\tilde O(d\sqrt T)$; its link assumption is stronger than
Li et al.'s Assumption~1; it uses an initial phase to make the design
invertible and a projection step.
\end{description}

\section*{Draft T2: a Lipschitz bound on scaled-tanh relative weights}

Write $\omega_\theta(z)=\operatorname{sech}^2(\phi(z)^\top\theta/\sqrt W)$ as in
\Cref{cor:tanh-relative-weights}.

\begin{lemma}[Relative-weight Lipschitz bound]\label{lem:tanh-weight-lipschitz}
For every $z$, $W>0$ and $\theta_1,\theta_2$,
\[
\frac{\omega_{\theta_1}(z)}{\omega_{\theta_2}(z)}
\le\exp\!\Big(\frac{2|\phi(z)^\top(\theta_1-\theta_2)|}{\sqrt W}\Big).
\]
If moreover $|\phi(z)^\top\theta_i|\le B_\phi R$ for $i=1,2$, the factor $2$
can be replaced by $2\tanh(B_\phi R/\sqrt W)$.
\end{lemma}

\begin{proof}[Proof {[checked]}]
Let $u_i=\phi(z)^\top\theta_i/\sqrt W$ and $h(u)=\log\operatorname{sech}^2u=-2\log\cosh u$.
Then $h'(u)=-2\tanh u$, so $|h'|\le2$ on $\mathbb R$ and $|h'|\le2\tanh x$ on
$[-x,x]$ because $\tanh$ is odd and increasing. By the mean value theorem,
$|h(u_1)-h(u_2)|\le\sup|h'|\,|u_1-u_2|$ over the segment between $u_1$ and
$u_2$, which lies in $[-x,x]$ when both endpoints do. Exponentiating
$h(u_1)-h(u_2)\le|h(u_1)-h(u_2)|$ gives both claims.
\end{proof}

The constant $2$ is asymptotically attained: for fixed $\epsilon>0$,
$\cosh^2(u)/\cosh^2(u-\epsilon)\to e^{2\epsilon}$ as $u\to\infty$ [checked].

\begin{corollary}[Lipschitz endpoint certificate]\label{cor:tanh-lipschitz-endpoint}
Under the hypotheses of \Cref{cor:tanh-relative-weights}, let
$M_t=\max_{s<t}|\phi(z_s)^\top(\btheta_t-\btheta_s)|$ ($M_1=0$) and
$D_t^{\rm lip}=4\tanh(B_\phi R/\sqrt W)\,M_t/\sqrt W$. Then
$e^{-D_t^{\rm lip}}\bVbar_t\preceq\bV_t\preceq e^{D_t^{\rm lip}}\bVbar_t$, so
$c_{\mp,t}=e^{\mp D_t^{\rm lip}}$ satisfy \eqref{eq:relative-endpoint-sandwich}
and $d_{\rm Th}(\bVbar_t,\bV_t)\le D_t^{\rm lip}\le4M_t/\sqrt W$.
\end{corollary}

\begin{proof}[Proof {[checked]}]
Put $c=e^{-D_t^{\rm lip}}\le1$. Then
$\bV_t-c\bVbar_t=\lambda(1-c)\bI+\sigma^{-2}\sum_{s<t}
(\omega_{\btheta_t}(z_s)^2-c\,\omega_{\btheta_s}(z_s)^2)\phi(z_s)\phi(z_s)^\top$.
The first term is PSD. By \Cref{lem:tanh-weight-lipschitz} with the bounded
index ($\|\phi\|\le B_\phi$, $\|\btheta\|\le R$),
$\omega_{\btheta_t}(z_s)^2/\omega_{\btheta_s}(z_s)^2\ge
e^{-4\tanh(B_\phi R/\sqrt W)|\phi(z_s)^\top(\btheta_t-\btheta_s)|/\sqrt W}\ge c$,
so every coefficient is nonnegative and each summand is PSD. The upper
inequality is symmetric with $c^{-1}\ge1$. The Thompson-distance statement is
its definition; the last inequality uses $\tanh\le1$.
\end{proof}

\paragraph{Comparison with \Cref{cor:tanh-relative-weights}.} Let
$x=B_\phi R/\sqrt W$.
\begin{itemize}
\item \emph{Uniform certificate} $D^{\rm unif}=4\log\cosh x$. $D_t^{\rm lip}<D^{\rm unif}$
iff $M_t<\sqrt W\log\cosh(x)/\tanh(x)$; with $B_\phi=R=1$ the threshold is
$0.505$, $0.570$, $0.687$ at $W=16,1,0.25$ [computed]. Since
$M_t\le2B_\phi R$, the worst case of $D^{\rm lip}$ is $8x\tanh x$, which exceeds
$D^{\rm unif}$ for every $x>0$ because $\log\cosh x\le x\tanh x<2x\tanh x$
[checked: $\frac{d}{dx}(x\tanh x-\log\cosh x)=x\operatorname{sech}^2x\ge0$].
So the uniform certificate is tighter for large index movement, and the
Lipschitz certificate is tighter when $M_t$ is below the threshold.
\item \emph{Category certificate} \eqref{eq:tanh-category-relative}. It is never
looser than $D^{\rm lip}$ [checked]: for an observed category,
$g_{t,c}/f_{t,c}=\omega_{\btheta_t}(c)^2\big/\big(n_{t,c}^{-1}\sum_{s<t:\phi(z_s)=\phi_c}\omega_{\btheta_s}(c)^2\big)$
divides by a mean of the $\omega_{\btheta_s}(c)^2$, which lies between their
minimum and maximum, so $g_{t,c}/f_{t,c}\in[e^{-D_t^{\rm lip}},e^{D_t^{\rm lip}}]$; adding $1$ to
the min and max keeps both factors in that interval. The Lipschitz
certificate does not need finite categories; it bounds the certificate by
index movement, and through $M_t\le B_\phi\max_{s<t}\|\btheta_t-\btheta_s\|_2$, by
parameter movement.
\item \emph{Remark (Hessian path certificate, \Cref{cor:hessian-path-new}).} With
$L_g=c_h/\sqrt W$, $c_h=4B_\phi^2/(3\sqrt3)$ and
$\max_{s<t}\|\btheta_t-\btheta_s\|\le\sqrt{Q_t}$,
$4M_t/\sqrt W\le4B_\phi\sqrt{Q_t}/\sqrt W=\frac{3\sqrt3\,\sigma\sqrt\lambda}{2B_\phi}\bar D_t$
[checked]. With $\sigma=0.25$, $\lambda=B_\phi=1$ the factor is $0.650$.
\end{itemize}

\section*{Draft T3: the uniform endpoint factor when $W\ge cT$}

\begin{corollary}\label{cor:tanh-endpoint-regime}
If $W\ge cT$ for some $c>0$, then
$\exp(4\log\cosh(B_\phi R/\sqrt W))=\cosh^4(B_\phi R/\sqrt W)\le\exp(2B_\phi^2R^2/(cT))$.
\end{corollary}

\begin{proof}[Proof {[checked]}]
For $y\ge0$ let $f(y)=y^2/2-\log\cosh y$; $f(0)=0$ and $f'(y)=y-\tanh y\ge0$,
so $\log\cosh y\le y^2/2$, and evenness extends this to all $y$. Hence
$4\log\cosh(B_\phi R/\sqrt W)\le2B_\phi^2R^2/W\le2B_\phi^2R^2/(cT)$; $\exp$ is
increasing.
\end{proof}

\paragraph{Implication under \Cref{cor:global-near-linear-new}'s assumptions.}
Take the scaled-tanh specialization ($c_\mu=c_g=c_h$, $G=B_\phi$), exact current
curvature, exact solves and exact maximization ($\kappa_{\pm,t}=\alpha_t=1$,
$\xi_t=0$), and score with \Cref{cor:relative-endpoint-ucb} using the uniform
certificate $c_{\pm,t}=\ell^{\mp2}$, $\ell=\operatorname{sech}^2(B_\phi R/\sqrt W)$.
Then $h_t=1+\ell^{-2}=1+\cosh^4(B_\phi R/\sqrt W)\le\bar h:=1+e^{2B_\phi^2R^2/(cT)}$.
Using \eqref{eq:relative-endpoint-regret} and
\Cref{prop:frozen-width-ucb} as stated in the paper [not checked: their proofs
were not re-read in this session]:
\[
R_T^{\rm endpoint}\le\bar h\sqrt{\Big(\sigma^2+\frac{G^2}{\lambda}\Big)\gamma_T\sum_t\beta_t^2}+2\sum_tb_t(a_t),
\qquad
R_T^{\rm frozen}\le2\sqrt{\Big(\sigma^2+\frac{G^2}{\lambda}\Big)\gamma_T\sum_t\beta_t^2}+2\sum_tb_t(a_t),
\]
each on its own policy's trajectory. So, as functions of the same $\gamma_T$,
$\beta_t$ and $b_t$, the transported exploration term is at most
$\bar h/2=\frac12(1+e^{2B_\phi^2R^2/(cT)})=1+O(1/(cT))$ times the frozen one;
the bias terms are identical. This compares bound expressions; the realized
$\gamma_T$, $\beta_t$, $b_t$ differ between the two trajectories.
Bounding $\gamma_{t-1}\le\Gamma_{T,d}$ (\eqref{eq:finite-d-gamma-new}),
$\epsilon_{\rm lin}(s),b_t\le2c_hR^2/\sqrt W$ (Taylor with $\|\btheta^\circ-\btheta_s\|\le2R$)
and $\epsilon_{\rm mis}=0$ in \eqref{eq:beta-corrected-new} gives [checked]
\[
R_T^{\rm endpoint}\le\bar h\Big[\sqrt{\Gamma_{T,d}+2\log(1/\delta)}+\sqrt\lambda R
+\frac{2c_hR^2\sqrt T}{\sigma\sqrt W}\Big]\sqrt{\sigma^2+\frac{G^2}{\lambda}}\sqrt{T\Gamma_{T,d}}
+\frac{4c_hR^2T}{\sqrt W}
\]
with probability at least $1-\delta$, which is \eqref{eq:global-near-linear-new}
with $2e^{D_{T,W}}$ replaced by $\bar h$. With $W=cT$, the paper's
$D_{T,W}=4c_hR\sqrt{T-1}/(\sigma\sqrt{\lambda cT})$ stays below
$4c_hR/(\sigma\sqrt{\lambda c})$, about $12.3/\sqrt c$ at the study's constants,
while $\bar h\to2$. At the study's $T=1000$ and $W=16,1,0.25$:
$D_{T,W}=97.3,389.3,778.6$ and $1+\cosh^4(1/\sqrt W)=2.13,6.67,201.3$ [computed].
Nothing is claimed here about analyses other than the paper's.

\section*{Draft T4: related work and positioning}

\paragraph{Related work (draft).}
Scaled tanh is a single-index model with a known increasing link, the form
studied by generalized linear bandits. Li et al.\ (2017) analyze UCB-GLM and
SupCB-GLM under conditionally sub-Gaussian noise, a link with
$\inf\dot\mu=\kappa>0$ near the true parameter, bounded first and second link
derivatives, i.i.d.\ contexts with a nondegenerate design, and an estimator
that solves the GLM score equation; their bounds scale as
$\tilde O(d\sqrt T/\kappa)$ and $\tilde O(\sqrt{dT\log K}/\kappa)$. For the logistic
link with Bernoulli rewards, Faury et al.\ (2020) replace the global $\kappa$ in
the leading term by local quantities using a Bernstein-type self-normalized
inequality, and Abeille et al.\ (2021) obtain a leading term
$d\sqrt{T/\kappa_\star}$ with a matching local lower bound. None of these
results applies as stated to our setting: the logistic analyses assume
Bernoulli rewards, and Li et al.\ state bounded rewards and a nondegenerate
context design that the study's 29-dimensional features do not have. Our
analysis also differs in its object: it allows an arbitrary predictable
parameter fitted with squared loss and corrects the prediction with a frozen
linear estimate, rather than analyzing a likelihood estimator.
% Filippi et al. (2010) to be added once the primary source is read.

\paragraph{Positioning (draft).}
Transport matters only when current and collection-time curvature differ. For
scaled tanh the two metrics differ only through squared link derivatives, and
a log-Lipschitz bound on those derivatives
(\Cref{lem:tanh-weight-lipschitz}) gives an endpoint certificate controlled by
how far the linear index has moved; the category certificate of
\Cref{cor:tanh-relative-weights} is never looser. In the near-linear regime
$W\ge cT$ that \Cref{cor:global-near-linear-new} targets, the uniform endpoint
factor is at most $e^{2B_\phi^2R^2/(cT)}$
(\Cref{cor:tanh-endpoint-regime}), so the endpoint-transported bound has the
frozen bound's form with the factor $2$ replaced by
$1+e^{2B_\phi^2R^2/(cT)}$, whereas the Hessian path certificate leaves the
factor $2e^{D_{T,W}}$ with $D_{T,W}=O(1/\sqrt c)$. Transport therefore costs at
most a $1+O(1/(cT))$ factor relative to frozen scoring under these
assumptions; it does not improve on frozen scoring in this analysis.
